\section{Counting and finding all relevant supports}\label{sec:enumeration}
The results in this section do not require quadratic functions or a graph.
Let $Z$ be a nonempty subset of $\{0,1\}^m$. The case $m=0$ consists of
one support and follows with the empty-product conventions.

\subsection{A near-optimal support count}
Assume the independent real random variables $\xi_i$ satisfy
\begin{equation}\label{eq:intervalmass}
 \Prob\{\xi_i\in J\}\le\phi |J|+\tau
 \quad\text{for every closed interval }J,
 \qquad \phi>0,\quad\tau\ge0.
\end{equation}
This allows atoms. For arbitrary deterministic costs $g_z\in\R$, put
\[
 F_z=g_z+\xi^Tz,\qquad v=\min_{z\in Z}F_z,\qquad
 A_a=\{z\in Z:F_z\le v+a\},\quad a\ge0.
\]

\begin{lemma}[Near-optimal count]\label{lem:count}
Under \eqref{eq:intervalmass},
\begin{equation}\label{eq:count}
 \E|A_a|\le(1+2\phi a+\tau)^m.
\end{equation}
All tied supports are counted separately.
\end{lemma}
\begin{proof}
A family $A\subseteq\{0,1\}^m$ \emph{shatters} $D\subseteq[m]$ if its
projections onto $D$ contain every binary pattern. Write $\Sh(A)$ for its
shattered coordinate sets. The classical sandwich inequality
\cite[Theorem~5]{KozmaMoran} gives
\begin{equation}\label{eq:sandwich}
 |A|\le|\Sh(A)|.
\end{equation}
For completeness, split $A$ on its last coordinate, obtaining projected
families $A_0,A_1$. Set $U=A_0\cup A_1$ and $V=A_0\cap A_1$. Then
$|A|=|U|+|V|$. Every set shattered by $U$ is shattered by $A$, and every
set shattered by $V$ extends, by the last coordinate, to a set shattered
by $A$. These two collections are disjoint. Induction proves
\eqref{eq:sandwich}, starting with the empty family, which shatters no set,
and the singleton family in dimension zero, which shatters the empty set.

We use the bit-class conditioning method of Beier and V\"ocking's generalized
isolating lemma \cite[conference version, Lemma~5]{BeierVocking}, combined
here with \eqref{eq:sandwich} and the interval-mass bound.
Fix $D\subseteq[m]$ and condition on $(\xi_i)_{i\notin D}$. For a pattern
$u\in\{0,1\}^D$, define
\[
 C_u=\min_{z\in Z:\ z_D=u}
          \left(g_z+\sum_{i\notin D}\xi_i z_i\right).
\]
If a class is empty, $D$ cannot be shattered. Otherwise these class costs
are finite and independent of $\xi_D$. If $A_a$ shatters $D$, every class
has a near-optimal representative, so its minimum lies in $[v,v+a]$.
Comparing the zero pattern with each unit pattern gives
\[
 |C_{e_i}+\xi_i-C_0|\le a\qquad(i\in D).
\]
These are fixed intervals of length $2a$ for the independent remaining
coordinates. Therefore
$\Prob\{D\in\Sh(A_a)\}\le(2\phi a+\tau)^{|D|}$.
Taking expectations in \eqref{eq:sandwich} and summing over all $D$ gives
\eqref{eq:count} by the binomial formula.
\end{proof}

If $\xi_i$ is uniform on the endpoint-inclusive grid
\begin{equation}\label{eq:noisegrid}
 \left\{-\sigma+\frac{2\sigma j}{N-1}:j=0,\ldots,N-1\right\},
 \qquad N\ge2,\quad\sigma>0,
\end{equation}
then
\begin{equation}\label{eq:gridmass}
 \phi=\frac1{2\sigma},\qquad\tau=\frac1N
\end{equation}
are valid. Indeed an interval of length $d$ contains at most
$d(N-1)/(2\sigma)+1$ grid points. This direct estimate includes grid
endpoints and exact ties.

\subsection{Certified enumeration with an additive oracle}
A partial assignment $p$ fixes some coordinates of $z$; let $Z_p$ be the
corresponding subset of $Z$. At accuracy $\varepsilon>0$, assume an oracle
either recognizes $Z_p=\varnothing$ or returns $z_p\in Z_p$ and its
\emph{exact} cost $V_p=F_{z_p}$ with
\begin{equation}\label{eq:oracle}
 \min_{z\in Z_p}F_z\le V_p\le\min_{z\in Z_p}F_z+\varepsilon.
\end{equation}
Thus $\ell_p=V_p-\varepsilon$ is a certified lower bound. The oracle must
work under arbitrary fixed-bit restrictions.

For a tolerance $\delta\ge0$, use the following procedure. Query the
unrestricted cell and keep its candidate value $U$ fixed. Store nonempty
cells in a heap keyed by $\ell_p$, caching their candidates. While the
smallest key is at most $U+\delta$, extract that cell and output its
candidate $z$. If its free coordinates in increasing order are
$j_1,\ldots,j_r$, create the $r$ children
\begin{equation}\label{eq:children}
 z_{j_h}=z^{\rm out}_{j_h}\ (h<\ell),\qquad
 z_{j_\ell}=1-z^{\rm out}_{j_\ell},\qquad \ell=1,\ldots,r,
\end{equation}
retaining the parent's fixed coordinates; $z^{\rm out}$ denotes the
extracted candidate. Query each child and insert it if nonempty. Stop when
the heap is empty or its smallest key exceeds $U+\delta$.
This is Lawler's first-difference partition \cite{Lawler}, with certified approximate
values used in the stopping test.

\begin{lemma}[Enumeration guarantee]\label{lem:enumeration}
The procedure terminates, has no repeated outputs, and returns a set $E$
satisfying
\begin{equation}\label{eq:inclusions}
 A_\delta\subseteq E\subseteq A_{\delta+2\varepsilon}.
\end{equation}
It uses at most $1+m|E|$ oracle calls.
\end{lemma}
\begin{proof}
The children in \eqref{eq:children} partition the parent cell except for
the extracted support: every other support has a unique first differing
free coordinate. Hence the stored cells always partition all unextracted
supports. Outputs are distinct, so the finite family guarantees termination.
Since $v\le U$, a remaining support of cost at most $v+\delta$ belongs to
a cell with lower bound at most its cost and hence at most $U+\delta$.
The stopping condition cannot hold while such a support remains.
Conversely an extracted candidate has
\[
 V_p=\ell_p+\varepsilon\le U+\delta+\varepsilon
                \le v+\delta+2\varepsilon.
\]
Each extraction creates at most $m$ child calls, proving the count.
\end{proof}

No exact ranking of the outputs is asserted. The guarantee holds even when
each oracle call returns the worst permitted approximate candidate. It is
pointwise in the entire noise vector; later probability bounds do not
condition on the adaptive search history.

\subsection{A complete dictionary on a parameter domain}
Let $T$ be a nonempty compact metric space and let deterministic functions
$q_z:T\to\R$ be $L$-Lipschitz for a common $L\ge0$. Define
\[
 F_z(t)=q_z(t)+\xi^Tz,\qquad V(t)=\min_{z\in Z}F_z(t).
\]
The set of supports optimal somewhere is measurable: for each $z$, its
activity is equivalent to
$\min_{t\in T}(F_z(t)-V(t))=0$, and this minimum is continuous in the
finite noise vector. For the expectation statements, the oracle's
selected support must be measurable in the noise; all remaining ties are
resolved by fixed deterministic orders.

\begin{theorem}[Complete dictionary from restricted approximation]
\label{thm:dictionary}
Suppose \eqref{eq:intervalmass} holds, $n\ge\max\{1,m\}$, and
$\tau\le1/(2n)$. Set
\begin{equation}\label{eq:accuracy}
 \varepsilon=\delta=\frac1{12n\phi}.
\end{equation}
For $L>0$, take a deterministic $\varepsilon/(2L)$-net
$t_1,\ldots,t_J$ in $T$. For $L=0$, take any one point. Run the preceding
enumeration at every $t_j$, using \eqref{eq:oracle} for $F_z(t_j)$, and let
$\mathcal D=\bigcup_j E_j$. Then, for every noise realization,
\begin{equation}\label{eq:dictionaryexact}
 \bigcup_{t\in T}\argmin_{z\in Z}F_z(t)\subseteq\mathcal D,
 \qquad V(t)=\min_{z\in\mathcal D}F_z(t)\quad(t\in T).
\end{equation}
Moreover,
\begin{equation}\label{eq:dictionarysize}
 \E\sum_{j=1}^J|E_j|\le eJ,
 \qquad\E[\text{number of oracle calls}]\le J(1+em).
\end{equation}
\end{theorem}
\begin{proof}
If $z$ minimizes at $t$ and $t_j$ is within $\varepsilon/(2L)$, compare it
with a minimizer $y$ at $t_j$:
\[
 F_z(t_j)\le F_z(t)+\varepsilon/2
           \le F_y(t)+\varepsilon/2
           \le F_y(t_j)+\varepsilon.
\]
Lemma~\ref{lem:enumeration} therefore outputs $z$ at $t_j$. For $L=0$
the same conclusion holds at the single point because every branch is
constant. This proves inclusion, even for supports optimal only at ties
or on the boundary. Every extra output is a valid member of $Z$, which
proves equality of envelopes. At each deterministic net point,
\[
 \E|E_j|\le\E|A_{3\varepsilon}(t_j)|
 \le(1+6\phi\varepsilon+\tau)^m
 \le(1+1/n)^m\le e.
\]
Sum these estimates and use the oracle-call bound. The same random vector
is used at every point; independence between points is unnecessary.
\end{proof}

For a rational box $[-R,R]^k$ with $k,R,L>0$, divide each coordinate into
$h=\max\{1,\lceil2kRL/\varepsilon\rceil\}$ equal intervals and use their
product midpoints. The Euclidean covering radius is at most
$\sqrt{k}R/h\le kR/h\le\varepsilon/(2L)$, and
\begin{equation}\label{eq:netsize}
 J=h^k\le(1+2kRL/\varepsilon)^k.
\end{equation}
If $k=0$ or $R=0$, the box is a singleton. Rational $R,L,\varepsilon$
give rational net points of polynomial encoding length.

To turn the theorem into an algorithm for functions, one also needs
efficient construction of each branch representation from its support.
Only a first-moment bound is required: at a single net point at most
$1+m2^m$ cells are generated, so every heap operation uses
$O(m+\log(m+1))$ comparisons, even on a large realization. Support strings
can be deduplicated in a binary trie in $O(m)$ time per output. Thus the
additional work is linear in the output count, up to deterministic
polynomial factors. If exact evaluation, branch construction, and each
restricted oracle call have uniformly polynomial bit cost, the expected
bit cost is polynomial whenever $J$ is polynomial. This conclusion does
not require higher moments or construction of optimality regions.
